% Introduction. Written after the results gate from the result claims; every number is a macro.
% Claims: r-sel-candidate, r-v-h2-headline, r-h2-emp-*, r-v-h2 (label probabilities), r-h90-h2,
% r-mde-h2, r-plc-nonff-sign, r-v-s1, r-v-s2, r-v-h1, r-h1-wage-*, r-v-h2-note, r-ph-h1-joint,
% r-v-h3, r-v-n-not-robust.
\section{Introduction}
\label{sec:intro}

On April 1, 2024 California raised the minimum wage of workers at fast-food chains with at least
60 establishments nationwide to \$20 an hour, under Assembly Bill 1228. The studies that
estimate its effect on employment from the same administrative data, the Quarterly Census of Employment and Wages
(QCEW), disagree. \citet{clemens2026employment} find that California's fast-food employment fell
by about 3.5 percent relative to the rest of the country over two years; \citet{dube2026labor}
finds an effect near zero with a different comparison group, outcome normalisation and
estimator. The disagreement has been decomposed into these choices
\citep{dube2026labor,clemens2026synthetic,sosinskiy2026effects}, which means that anyone who
re-estimates the effect knows in advance which choices give which answer.

We ask whether AB 1228 changed employment in California's private limited-service restaurants
(NAICS 722513) over its first eight quarters, 2024Q2--2026Q1, relative to the same industry in
other states, and whether the effect changed between the first and the second year, when the
law's premium over California's general minimum wage fell. Because the specification choices
and their published consequences were known, we did not choose a specification. Before any
outcome after 2023Q3 was computed we fixed the estimand, the base quarter, decision rules for
every label, an inference procedure built for one treated state, and a rule that picks the
primary estimator from 24 candidates covering both published approaches by how well each
reproduced California's own path at placebo dates in 2016--2017, using only data through 2019.
Both published approaches are also re-estimated as named secondary analyses. The study is
pre-specified but not blind: its authors knew the published estimates and a public descriptive
series through March 2026 before the design was fixed (Section~\ref{sec:prespec}). The data add
one preliminary quarter, 2026Q1, to the latest published window.

The answer is that this design cannot tell. The rule selected \cand{DDT-POPW-E-R50}, a
difference in differences with a linear trend in the gap, population-weighted donors and all
other states and the District of Columbia as the comparison pool. Its pooled effect on log
employment is $\HTwoAtt$ log points ($\HTwoPct$ percent, about $\HTwoJobs$ jobs), with a 95
percent interval from $\HTwoLoNinetyFive$ to $\HTwoHiNinetyFive$ ($\HTwoPctLo$ to $\HTwoPctHi$
percent). Under the pre-specified rule the label is
``inconclusive''. Employment outside limited-service restaurants, by contrast, ran above its
counterfactual, by $\PlcNonffAtt$ log points with an interval that excludes zero, so the label
is qualified as not separable from a California-wide shock; a shared deviation of that sign
would bias the employment estimate toward zero or upward, not toward a decline. The two named
secondary analyses, the detrended comparison of Clemens and coauthors and the synthetic
difference in differences of Dube, both re-estimated with our timing and inference, are also
``inconclusive'' ($\SOneAtt$ and $\STwoAtt$), and no robustness check changes the sign category.

The equivalence test against a 2 percent bound was meant to separate the two published
readings, and it could not. The power run, made on data through 2019 before any post-period
outcome was read, put the half-width of the 90 percent interval at $\HNinetyHTwo$ log points,
above the bound of 0.02, so the label ``no effect as large as 2\%'' had probability
$\ProbTauZeroUndSmaller$ whatever the true effect; ``inconclusive'' had probability
$\ProbTauZeroUndUndet$ with no effect and $\ProbTauAltUndUndet$ with a decline of 3.5 percent.
The label is therefore not evidence against either published estimate, both of which lie inside
our interval. The width comes from the design's inference more than from the data: the standard
error is the larger of a population-scaled and an unscaled placebo standard error, chosen to keep
the level of the tests for California, the largest state, and at the placebo dates California's
own six $t$ statistics had a root mean square of $\KDiagHTwo$ with it, against
$\KDiagHTwoScaled$ with the scaled one. With the scaled standard error alone, which the design
does not use because it over-rejects if shocks have a state-level component, the interval would
lie entirely below zero; the pre-2020 diagnostics are mixed on whether such a component is present
(Sections~\ref{sec:res-inference} and~\ref{sec:res-h2}).

The first stage is not established either. The pooled effect on log average weekly earnings in
722513 is $\HOneAtt$ log points ($\HOnePct$ percent, about \$$\HOneDollars$ a week), with
$t = \HOneT$. That is below the pre-specified one-sided critical value $\HOneCrit$, which the
placebo calibration set far above $t$ at the estimated degrees of freedom ($\TnuHOneNinetyFive$),
so hypothesis H1 is not supported, and by the design's rule the employment result is reported as
uninformative about the wage floor. This is not a finding that weekly earnings did not rise. Two
post hoc checks bear on the verdict, and they point in opposite directions. Removing one
comparison state's placebo tests, whose trend windows straddle a 2015 drop in that state's
weekly-wage series, lowers the critical value to $t$ at the degrees of freedom, and the test
would then reject; but that state was singled out after the results were seen. The selected
estimator's wage gap drifted before 2020 by $\PhDriftWageSlope$ log points a quarter, which
extrapolated over the post window is $\PhDriftWageImpliedPost$ against the estimate of
$\HOneAtt$; but a linear extrapolation is not a test, and the pre-2020 years contain
California's own minimum-wage ramp. Neither check changes the pre-specified verdict
(Section~\ref{sec:res-posthoc}). The
comparison of year two with year one is uninformative by the pre-specified power rule.

The contribution is therefore modest and mainly about what a state-level comparison with this
design's conservative inference for one treated state can and cannot settle.
We add an estimate with a rule-selected primary specification, inference whose level was checked
on placebo states before the post-period data were read, and one more quarter of data; and we
document, before seeing the results, that with this inference the precision is too low to
separate a decline of about 3.5 percent from no effect, and after them, how much of that comes
from the choice of standard error. The
published decompositions of the disagreement are cited, not repeated
(Section~\ref{sec:literature}).

Section~\ref{sec:background} describes the law and its coverage, Section~\ref{sec:literature} the
literature, Section~\ref{sec:data} the data and Section~\ref{sec:strategy} the pre-specified
strategy. Section~\ref{sec:results} reports the selection, the inference and power checks and the
three hypotheses; Section~\ref{sec:robustness} the secondary analyses, robustness and placebo
checks; Section~\ref{sec:conclusion} concludes, and the last section lists the limitations and
every deviation from the approved design.
